% Chapter 1
\chapter{Introduction}
\label{chap:introduction}

This chapter establishes the foundation for investigating observability frameworks in containerized \ac{CICD} environments. It frames the research problem within both theoretical and organizational contexts, articulates the specific challenges arising from heterogeneous platform deployments, and defines the scope and objectives that guide this work.

The chapter begins by presenting the operational challenges that motivate this research, then examines the design trade-offs and architectural decisions required when implementing comprehensive observability for on-premise containerized infrastructure. Research questions, intended contributions, and methodological approach are subsequently outlined, providing a roadmap for the investigation documented in this preparation report.

\section{Introduction and Problem Framing}

Infrastructure failures in containerized \ac{CICD} platforms create immediate operational challenges: determining which service failed, identifying root cause, and restoring functionality before development teams are blocked. In organizations running on-premise \ac{CICD} infrastructure, diagnostic processes often require manual log inspection across multiple servers, timestamp correlation between disconnected systems, and inference from incomplete information. This fragmentation extends detection and resolution times, directly affecting development velocity and release schedules.

While cloud-native observability solutions provide comprehensive telemetry collection, on-premise deployments face distinct challenges: heterogeneous platform requirements (mixed Linux and Windows hosts), integration complexity across multiple specialized services, and the need to balance observability depth with operational overhead. Organizations with Docker access have multiple instrumentation options available, ranging from lightweight container-level monitoring to comprehensive host-level instrumentation, each presenting different trade-offs in terms of implementation complexity, resource consumption, and diagnostic capability.

This project investigates observability approaches for containerized \ac{CICD} platforms deployed across heterogeneous infrastructure. The work examines architectural patterns for collecting, correlating and analyzing telemetry from multiple specialized services, evaluating different instrumentation strategies and their practical implications. The investigation focuses on a four-server on-premise deployment running mixed Linux and Windows hosts, where operations staff now have \texttt{docker} group access on the Linux hosts and must design an observability solution that balances comprehensive visibility with operational maintainability.

The central research problem is: \textit{How can an integrated observability framework be designed and evaluated to improve reliability and reduce incident response times in heterogeneous containerized \ac{CICD} environments, and what are the trade-offs between different instrumentation approaches?}


\section{Context and Motivation}

The motivation for this research emerges from the intersection of theoretical understanding of observability frameworks and the operational realities encountered in specific organizational deployments. This section examines both the conceptual foundations established in software engineering research and the practical constraints that shape observability requirements in the target environment. The two subsections establish the theoretical basis for observability in \ac{CICD} systems and describe the specific organizational context that defines the problem space for this investigation.

\subsection{Theoretical Context}
Continuous integration and continuous delivery (\ac{CICD}) pipelines constitute critical infrastructure for software development organizations. Research in software engineering has established that pipeline reliability and observability directly influence development velocity, defect detection rates and time-to-market~\cite{forsgren2024dora}. Modern \ac{CICD} platforms typically comprise multiple specialized services: source control management (such as GitLab), build automation (such as Jenkins), static analysis (such as SonarQube), artifact management (such as Nexus) and security scanning (such as Dependency-Track). Each service generates distinct telemetry streams (metrics, logs, traces) that must be collected, correlated and analyzed to maintain operational visibility.

Observability frameworks follow established patterns: metrics for quantitative monitoring, structured logging for event analysis, and distributed tracing for transaction flow reconstruction. Cloud-native observability assumes elastic infrastructure, privileged host access and homogeneous execution environments, enabling instrumentation patterns such as node exporters with host access, kernel-level tracing and infrastructure-as-code deployment.

\subsection{Organizational Reality}
The organization where this project is developed manages on-premise \ac{CICD} infrastructure serving multiple development teams across different business units. The platform runs on four physical servers: three Linux hosts and one Windows Server 2019 host. All \ac{CICD} services operate as Docker~\cite{docker2014} containers to facilitate deployment and isolation.

For a substantial period, operations staff did not have Unix \texttt{docker} group membership on the Linux hosts. Without that membership, routine use of the Docker API for container inspection, log retrieval, and co-located deployment of observability tooling was impractical, which tightened the effective access posture even though services themselves already ran in containers. This situation has since changed: the operations team now has \texttt{docker} group access on all Linux hosts, relaxing several practical constraints and enabling container lifecycle management, deployment of monitoring stacks as containers, and---where organizational policy permits---exploration of more intrusive instrumentation patterns (for example privileged monitoring containers or host-mounted volumes) without assuming a single omnipotent administrative account across the whole estate.

The environment nevertheless remains constrained relative to idealized cloud-native assumptions: mixed Linux and Windows workers, on-premise operation without tenant-wide orchestrator administration comparable to fully managed Kubernetes, and change-management boundaries that still limit what can be deployed in production-like paths. Together with multiple specialized services (GitLab, Jenkins, SonarQube, Nexus, Dependency-Track), these factors create architectural complexity for telemetry collection, correlation, and storage strategies.

Currently, the organization still relies predominantly on ad hoc diagnostics: application-level logs accessed through the Docker CLI where needed, manual inspection of container resource usage, and reactive troubleshooting after user-reported incidents. Even with improved Docker API access, there is not yet an integrated observability plane (unified ingestion, retention, alerting, and cross-signal navigation). This reactive posture results in prolonged incident detection (hours rather than minutes) and diagnosis cycles requiring coordinated investigation across disconnected log sources, and it limits proactive capacity planning.

The practical impact manifests in three ways: delayed detection of service degradation until user complaints escalate, extended root cause analysis requiring manual correlation of logs from five or more services, and inability to establish proactive capacity planning due to lack of historical resource utilization data. Development teams experience unpredictable pipeline failures attributed to infrastructure issues that operations teams cannot rapidly diagnose or prevent. The organization can therefore pursue a fuller range of container-centric observability designs than under the earlier access posture, but must still evaluate architectural approaches against operational risk, maintainability, and the remaining platform constraints.

\section{Problem Statement}

The problem addressed in this research comprises interrelated technical and organizational dimensions that collectively define the observability challenge in constrained \ac{CICD} environments. This section structures the problem statement across three complementary perspectives: the technical instrumentation challenges arising from administrative boundaries and platform heterogeneity (including a historically tighter access posture that has been partially relaxed through \texttt{docker} group membership for operations staff on Linux hosts), the organizational impact manifested through operational incidents and diagnostic delays, and the explicit boundaries that define what falls within and outside the scope of this investigation.

\subsection{Technical Problem}
Containerized \ac{CICD} platforms deployed on-premise with heterogeneous infrastructure present observability challenges requiring systematic investigation and design. With Docker group access for operations on the Linux hosts, multiple instrumentation approaches become feasible, each presenting distinct trade-offs. The technical problem comprises three interconnected dimensions:

\begin{enumerate}
    \item \textbf{Architectural design decisions:} Multiple patterns exist for telemetry collection: sidecar containers, privileged monitoring containers with host access, external agents, or hybrid approaches. Each pattern presents trade-offs between implementation complexity, resource overhead, maintainability, and diagnostic capability. Selecting and justifying appropriate patterns requires systematic evaluation of alternatives against defined criteria.
    \item \textbf{Telemetry integration and correlation:} Each \ac{CICD} service generates independent log streams in different formats. Designing centralized collection, storage and correlation mechanisms requires decisions about data flow architecture, schema normalization, retention policies, and query interfaces. Integration patterns must accommodate service-specific characteristics while enabling cross-service correlation for incident diagnosis.
    \item \textbf{Cross-platform heterogeneity:} Linux and Windows containers exhibit different runtime APIs, metric exposure mechanisms and log collection patterns. Architectural decisions must address platform-specific differences while maintaining unified telemetry schemas and operational procedures. The solution must provide consistent observability across platforms without requiring separate implementations for each environment.
\end{enumerate}

\subsection{Organizational Problem}
The organization experiences concrete operational difficulties from limited observability. Recent incidents illustrate the problem severity:

\begin{itemize}
    \item A GitLab database connection pool exhaustion incident in November 2025 was detected 3.5 hours after initial symptoms when multiple development teams reported failed pipeline executions. Root cause analysis required manual correlation of GitLab application logs, PostgreSQL (a relational database management system) container logs and Docker runtime inspection across two hosts.
    \item Intermittent Jenkins build failures in October 2025 attributed to insufficient executor capacity were diagnosed through trial-and-error resource adjustments over multiple days. Historical resource utilization data was unavailable to inform capacity planning decisions.
    \item A disk space saturation event in September 2025 affecting the Nexus artifact repository caused cascading build failures. Detection occurred only after builds began failing with cryptic upload errors. The absence of proactive disk usage monitoring prevented preventive action.
\end{itemize}

These incidents share common characteristics: detection relies on user reports rather than automated alerting, diagnosis requires time-consuming manual investigation without correlation tools, and preventive measures are hindered by lack of historical operational data. The organization estimates average \ac{MTTD} of 2--4 hours and \ac{MTTR} of 4--8 hours for infrastructure-related incidents.


\subsection{Scope and Boundaries}
This research is scoped to containerized environments with mixed operating systems (Linux and Windows) deployed on-premise. The investigation examines observability architectural patterns and evaluates different instrumentation approaches for their effectiveness, resource efficiency, and operational maintainability. The work focuses on \ac{CICD} platform services (source control, build automation, static analysis) rather than general-purpose application monitoring. Evaluation will be conducted within a single organizational context; generalizability to other platform configurations and organizational settings remains to be established through future research. The research does not address cloud-managed observability platforms, as the focus is on on-premise deployments where architectural decisions and trade-offs must be explicitly addressed.

\section{Intended Contributions}
The primary contributions planned are:
\begin{itemize}
    \item A reference architecture specifying patterns for collecting metrics, logs and traces from containerized services in mixed OS on-premise environments, including comparative analysis of alternative instrumentation approaches (sidecar, privileged containers, host-level agents) with documented trade-offs and selection criteria.
    \item A reproducible prototype demonstrating the selected architecture, delivered as container images, deployment templates and operational documentation, with procedures for deployment and troubleshooting.
    \item An empirical evaluation measuring changes in \ac{MTTD}, \ac{MTTR}, pipeline throughput and instrumentation overhead, including documented measurement procedures for reproducibility and assessment of different architectural alternatives where feasible.
    \item Practical guidance and configuration artifacts to support adoption by operations teams managing on-premise \ac{CICD} infrastructure, including decision frameworks for selecting appropriate observability patterns based on operational context, validated through practitioner feedback.
\end{itemize}

\section{Research Objectives}
The objectives guiding this project are:
\begin{itemize}
    \item \textbf{O1:} Design and evaluate a reference architecture for containerized \ac{CICD} observability, including systematic comparison of alternative instrumentation patterns (sidecar, privileged containers, host-level monitoring) with documented trade-off analysis and selection criteria.
    \item \textbf{O2:} Develop a reproducible prototype of container images and deployment templates instrumenting representative \ac{CICD} services, with documented configuration procedures and deployment automation.
    \item \textbf{O3:} Design and execute empirical evaluation comparing \ac{MTTD}, \ac{MTTR}, telemetry pipeline throughput and instrumentation overhead. Document measurement procedures, baseline collection methods and fault injection protocols for reproducibility.
    \item \textbf{O4:} Characterize resource overhead through CPU, memory and network utilization measurement. Document trade-offs between telemetry granularity and resource consumption, assessing whether instrumentation overhead can be maintained below 5\%.
    \item \textbf{O5:} Provide configuration templates, deployment scripts and documented procedures for on-premise environments, including decision criteria for architectural pattern selection and tool configuration, validated through practitioner feedback.
\end{itemize}

Each objective is defined with measurable outcomes and explicit deliverables. Objective O3 includes concrete measurement procedures (baseline collection, controlled fault injection, paired comparisons) to support quantitative assessment. Objective O4 establishes an explicit threshold (5\% overhead target) to guide evaluation and inform adoption decisions; the research will assess whether this threshold can be achieved in practice.

\section{Research Questions}
This project addresses the following questions:
\begin{enumerate}
    \item\label{rq:design-tradeoffs} How can an integrated observability framework be designed for heterogeneous containerized \ac{CICD} environments, and what are the trade-offs between different instrumentation approaches (sidecar containers, privileged monitoring, host-level agents)?
    \item\label{rq:correlation-maintainability} Which architectural patterns and deployment strategies enable effective cross-server telemetry correlation in mixed OS environments while maintaining acceptable resource overhead and operational maintainability?
    \item\label{rq:mttd-mtt-baseline} What measurable improvements in \ac{MTTD} and \ac{MTTR} can be achieved through integrated observability compared to baseline reactive monitoring, and what are the resource costs and operational complexities introduced?
\end{enumerate}

\section{Methodology Overview}
This project follows a design science research approach~\cite{hevner2004design}, combining systematic literature review, comparative technology assessment, artifact design, prototype implementation and empirical evaluation. Design science is appropriate because it addresses a practical problem through creation and evaluation of a novel artifact (the observability framework) in a specific organizational context. The research integrates qualitative methods (literature synthesis, practitioner interviews) with quantitative methods (performance measurement, resource overhead characterization). The work is organized into five phases:

\begin{description}
    \item[Phase 1 (Problem Analysis and Literature Review):] A PRISMA-guided systematic review examines observability frameworks, containerized monitoring and \ac{CICD} reliability literature to identify existing approaches, documented gaps and constraint patterns relevant to operationally constrained deployments (including privilege-limited patterns emphasized in prior work), interpreted alongside the evolving access posture of the target organization.
    \item[Phase 2 (Technology Assessment):] Evaluation of candidate tools (including Prometheus~\cite{prometheus2016} for metrics, Grafana Loki for logs, Jaeger for distributed tracing, and \gls{OpenTelemetry} for instrumentation) against defined criteria: feasibility under available operational access (including Docker Engine access for the operations account on managed hosts, without assuming unrestricted root or full Kubernetes cluster administration unless explicitly approved), cross-platform support (Linux and Windows), resource overhead characteristics and integration complexity. Assessment produces a justified tool selection with documented trade-off analysis.
    \item[Phase 3 (Architecture Design):] Specification of component placement (exporters, collectors, storage, visualization), data flow patterns and correlation mechanisms. Deliverables include deployment diagrams, configuration templates and documented design decisions with rationale.
    \item[Phase 4 (Prototype Implementation):] Implementation and documentation of a reproducible prototype using selected instrumentation patterns. Deliverables include deployment automation scripts, configuration artifacts and operational documentation.
    \item[Phase 5 (Evaluation and Validation):] Quantitative evaluation of \ac{MTTD}, \ac{MTTR} and instrumentation overhead through controlled fault-injection experiments with documented baseline measurements. Qualitative validation through structured practitioner interviews captures operational acceptability and deployment constraints.
\end{description}

The methodology balances artifact creation (reference architecture and prototype) with empirical validation (quantitative metrics and qualitative feedback). Quantitative evaluation establishes measurable outcomes; qualitative data addresses operational acceptability and deployment considerations that quantitative metrics alone cannot capture.

\section{Ethical Considerations}
This work adheres to the ethical guidelines established by Instituto Superior de Engenharia do Porto (\ac{ISEP}) and the IEEE Code of Ethics \parencite{IEEEXploreFullText}. The research, analysis, experimental design, and all scientific contributions presented in this document are authored by the student. The student assumes full responsibility for the integrity, accuracy and validity of all research findings, interpretations and conclusions.

This research involves human subjects activities (structured practitioner interviews) and access to operational telemetry data. All activities will follow institutional ethical review procedures. Ethical clearance will be obtained before commencing interviews or data collection. Participants will provide documented informed consent and will be informed of their right to withdraw. Interview data and operational telemetry will be anonymized to remove identifying information. Data handling will conform to institutional policies and applicable data protection regulations, including GDPR where relevant.

Production deployments will be preceded by staging environment validation and change management approval to minimize operational disruption risk. The project observes data minimization principles: telemetry collection is limited to technical metrics necessary for evaluation objectives. Personal data collection, if required, will be explicitly justified, documented and subject to institutional review board approval.

When presenting empirical results, host identifiers, network topology details and sensitive configuration information will be anonymized to protect organizational confidentiality. Access to production systems and telemetry data will be governed by organizational access control policies. All collected data will be stored securely and retained only for the duration necessary to complete the research, after which it will be securely disposed of in accordance with institutional data retention policies.

\subsection{AI Usage Disclosure}
Artificial intelligence tools were employed in a limited supporting capacity during the preparation of this work. AI assistance was restricted to: generating search keyword suggestions to support systematic literature identification; assisting with title and abstract screening during PRISMA-guided literature review; offering alternative phrasings for sentence reformulation to improve clarity and readability; and performing spell checking and grammar verification. AI tools were not used to generate scientific arguments, formulate research contributions, design experiments, analyze results, or draw conclusions. All intellectual content, methodological decisions, and scientific judgments remain the original work of the author.

\section{Document Structure}
This preparation report documents the research plan, including problem framing, current-state diagnosis, literature review, methodology, and the planned design-to-evaluation path. The remainder of this document is organized as follows: Chapter~\ref{chap:current_state} characterizes the current operational state of the target environment; Chapter~\ref{chap:literature} presents the systematic literature review following \ac{PRISMA} methodology; Chapter~\ref{chap:planning} details methodology and requirements planning; Chapter~\ref{chap:design} presents tool selection and architecture design; Chapter~\ref{chap:implementation} specifies implementation strategy; Chapter~\ref{chap:evaluation} defines evaluation and a-priori metrics; and Chapter~\ref{chap:conclusion} summarizes the preparation phase and next steps.